\section{DFA: Definisi, Tabel Transisi, dan Simulasi}
\label{sec:ch06-dfa}

\subsection{Definisi}

\begin{definisi}[Deterministic Finite Automaton]
DFA adalah lima-tupel
\[
M = (Q, \Sigma, \delta, q_0, F),
\]
dengan \(Q\) himpunan state berhingga, \(\Sigma\) alfabet berhingga,
\(\delta : Q \times \Sigma \to Q\) fungsi transisi total, \(q_0 \in Q\)
state awal, dan \(F \subseteq Q\) himpunan state penerima.
\end{definisi}

Sifat deterministik berarti untuk setiap pasangan state dan simbol
terdapat tepat satu state tujuan. DFA tidak memiliki transisi
\(\varepsilon\) dan tidak memiliki pilihan ganda. Karena itu, DFA
mudah disimulasikan: cukup menyimpan satu state saat ini.

\subsection{Contoh: Bilangan \(a\) Genap}

Misalkan \(L = \{w \in \{a,b\}^{*} \mid w \text{ memuat jumlah } a
\text{ genap}\}\). DFA untuk bahasa ini memiliki dua state:
\(q_0\) untuk jumlah \(a\) genap dan \(q_1\) untuk jumlah \(a\)
ganjil. State awal dan penerima adalah \(q_0\).

\begin{figure}[H]
\centering
\begin{tikzpicture}[node distance=24mm]
  \node[accepting, initial] (q0) {\(q_0\)};
  \node[state, right=of q0] (q1) {\(q_1\)};
  \draw[arrow] (q0) edge[loop above] node[font=\small] {\(b\)} (q0);
  \draw[arrow] (q1) edge[loop above] node[font=\small] {\(b\)} (q1);
  \draw[arrow] (q0) -- node[above, font=\small] {\(a\)} (q1);
  \draw[arrow] (q1) -- node[below, font=\small] {\(a\)} (q0);
\end{tikzpicture}
\caption{DFA untuk bilangan \(a\) genap}
\label{fig:ch06-dfa-genap}
\end{figure}

\begin{table}[H]
\centering
\caption{Tabel transisi DFA pada \figref{fig:ch06-dfa-genap}}
\label{tab:ch06-dfa-genap}
\begin{tabular}{@{}c cc@{}}
\toprule
\textbf{State} & \textbf{\(a\)} & \textbf{\(b\)} \\
\midrule
\(\to *q_0\) & \(q_1\) & \(q_0\) \\
\(q_1\) & \(q_0\) & \(q_1\) \\
\bottomrule
\end{tabular}
\end{table}

\subsection{Fungsi Transisi Diperluas}

Fungsi transisi \(\delta\) didefinisikan untuk satu simbol. Kita
memperluasnya menjadi \(\hat{\delta} : Q \times \Sigma^{*} \to Q\)
dengan
\[
\hat{\delta}(q, \varepsilon) = q, \qquad
\hat{\delta}(q, wa) = \delta(\hat{\delta}(q,w), a).
\]
String \(w\) diterima jika \(\hat{\delta}(q_0, w) \in F\).

\subsection{Simulasi}

Simulasi DFA sangat sederhana:
\begin{enumerate}[label=\arabic*.,leftmargin=*]
  \item mulai dari \(q_0\);
  \item untuk setiap simbol masukan, pindah ke
        \(\delta(q, a)\);
  \item jika transisi tidak terdefinisi, tolak;
  \item setelah masukan habis, terima jika state berada di \(F\).
\end{enumerate}

Kompleksitas waktu simulasi adalah \(O(n)\), dengan \(n\) panjang
masukan, dan kompleksitas ruang \(O(1)\) selain tabel transisi.
Sifat inilah yang membuat DFA cocok untuk lexer berkecepatan tinggi.
